\section{Conclusion}\label{sec:conclusion}
We asked when an automated vulnerability detector can be trusted to clear code. Framing clearance as risk control gives a precise answer: it can be trusted to the extent that the miss rate was calibrated on data that resembles the code it will see. On exchangeable splits conformal clearance delivers its guarantee on real detector scores, clearing 54--55\% of functions at a 10\% miss budget. Under project, temporal and cross-dataset shift the same procedure misses $2.7\times$ the nominal share on average, covariate-shift weighting does not repair it, and about 100 labelled target positives do, at roughly half of the automation; adaptive conformal inference achieves the same long-run control on drifting streams from delayed feedback alone. The guarantee is also uneven across function sizes and vulnerability classes, sensitive to cosmetic code changes, and economically attractive only when a miss costs tens of reviews. Frozen CodeBERT scorers replicate these findings, and expose a formatting shortcut in Big-Vul that makes raw-text models look valid and efficient within the dataset and unreliable beyond it. Together with a leakage audit that shows cross-dataset performance can be overstated by up to a factor of six, these results argue for reporting miss-versus-clearance curves on shift-aware splits and for treating the miss budget as a monitored, re-calibrated operational quantity.

Future work will repeat the study with fine-tuned code language models and large language models (we could only use frozen CodeBERT embeddings), evaluate adaptive (online) conformal variants and better density-ratio estimators for shift, extend the analysis to other C/C++-external languages and to repository-level detection, and apply the same risk-control formulation to other automation decisions in the software supply chain.
